\section{Introduction}
\begin{figure*}[t]
  \centering
  \includegraphics[width=\textwidth]{overview-ral.pdf}
  \caption{\textbf{Predict, anchor, refine.}
  Real TUM frames enter a shared predictor. Its geometric outputs support
  tracking; its Gaussian outputs form local maps. A keyframe anchor carries
  each local map and its supervision cameras, so both follow pose corrections.
  The refiner samples tracked images, renders the assembled map, and updates
  Gaussian appearance. Section and equation references locate each step.
  Insets show saved maps before and after refinement at the same camera.
  Images are raster data; the diagram and labels are editable in PowerPoint.}
  \label{fig:overview}
\end{figure*}
A robot revisiting a room needs a map that preserves both its layout and its
appearance. Visual SLAM estimates the camera motion and scene geometry;
3D Gaussian splatting~\cite{kerbl20233dgs} adds a representation that can
render the scene from another viewpoint. Gaussian SLAM
systems~\cite{matsuki2024monogs,huang2024photoslam} build such maps by
optimising primitives as images arrive.

Feed-forward reconstruction offers a faster starting point.
Splatt3R~\cite{smart2024splatt3r} predicts dense Gaussians from two images
using MASt3R features~\cite{leroy2024mast3r}. Repeating this operation along
a video produces many local reconstructions. It does not yet produce one
consistent map. Different pairs can place the same surface at different
depths or assign it different colours. A loop closure adds a second problem:
correcting old camera poses also changes where their predictions belong.

Our central idea is to give a local map and its observations the same anchor.
We attach predicted Gaussians to keyframes in
MASt3R-SLAM~\cite{murai2025mast3rslam}. We also store each supervision camera
relative to a keyframe. When the pose graph moves that anchor, it moves both
the local scene and the cameras that observe it. Their relative transform is
preserved. This lets a separate refiner improve appearance while tracking
continues to estimate motion.

Splatt3R-SLAM follows three steps: predict local Gaussians, anchor them to
the pose graph, and refine the assembled map (Fig.~\ref{fig:overview}).
The main contribution is this shared-anchor mapping and supervision design.
We then study two ways to improve its starting point: adapting only the
Gaussian head and reducing opacity when predictions enter the map.
Finally, controlled component studies and common-camera comparisons test
whether the design preserves tracking, improves rendering, and delivers
useful results within an online budget. On Replica, the measured PSNR gain
over Photo-SLAM is 3.54\,dB; the dense maps also reveal a substantial memory
cost. These outcomes motivate a focus on map quality per primitive.

\section{Related Work}
\subsection{Learned geometry for SLAM}
DUSt3R~\cite{wang2024dust3r} predicts dense pointmaps from image pairs, and
MASt3R~\cite{leroy2024mast3r} adds matching features. MASt3R-SLAM uses these
priors for tracking, keyframe fusion, and a similarity-transform pose graph.
VGGT~\cite{wang2025vggt} extends joint prediction to multiple images, and
VGGT-SLAM~\cite{maggio2025vggtslam} incorporates it into SLAM.
These methods supply geometric structure. Splatt3R adds Gaussian appearance
to pairwise prediction. Our work connects that appearance model to the
changing keyframe poses of a complete sequence. We use MASt3R-SLAM for
tracking and evaluate trajectory accuracy as a control.

\subsection{Gaussian maps and appearance refinement}
MonoGS~\cite{matsuki2024monogs} estimates motion and appearance with Gaussian
splatting. Photo-SLAM~\cite{huang2024photoslam} combines an ORB-SLAM3-based
geometric system~\cite{campos2021orbslam3} with photorealistic mapping.
SplaTAM~\cite{keetha2024splatam} and Gaussian-SLAM~\cite{yugay2023gaussianslam}
study Gaussian reconstruction with RGB-D input. These systems allocate and
optimise primitives during mapping. Our map starts with dense predictions
and refines their parameters after assembly. This reduces the need to grow
geometry from sparse seeds, but retains many overlapping primitives.

Unblur-SLAM~\cite{zhang2026unblur} combines learned deblurring with
rendering-based test-time refinement. It addresses degraded input images;
we address the placement and refinement of pairwise predictions under pose
corrections. The shared-anchor construction is the link between the learned
initial estimate and sequence-level optimisation in our system.
